\chapter{State of the Art}\label{ch:literature}

An incident-response decision requires knowledge that applies to the affected environment, supports a defensible course of action, and remains available for examination. We review organisational knowledge sharing, machine-learning representations, retrieval-augmented generation, and evaluation from this perspective. The chapter offers a focused narrative synthesis rather than an exhaustive systematic review. Its comparison of cybersecurity studies considers the task, knowledge sources, model adaptation, and evaluation method, since improvements on different tasks cannot be interpreted as a common measure of incident-response effectiveness.

\section{Incident response and organisational knowledge}
NIST Special Publication~800-61r3 places incident response within continuing cybersecurity risk management. In its Cybersecurity Framework (CSF)~2.0 model, Detect, Respond, and Recover cover incident handling, while Govern, Identify, and Protect support preparation and risk reduction. Lessons learned inform improvements across these activities rather than belonging only to a final reporting stage. NIST also distinguishes an observable event from an incident: interpreting evidence is necessary to establish whether an incident has occurred~\cite{nist2025incident}. This distinction matters for knowledge support. Retrieving a description of an attack does not establish that the attack explains an observed event.

Consider an analyst investigating suspected account compromise. A threat report can explain a technique, an internal procedure can specify an approved response, and a previous investigation can document exceptions encountered in a similar environment. The analyst still needs to assess which observations support the suspicion and whether the procedure applies. This example separates case evidence from reusable knowledge. Event records describe what was observed; response knowledge supplies the concepts, procedures, and prior experience needed to interpret those records. Their connection is a judgement about applicability, not simply a match between keywords.

We use \emph{knowledge-management system} (KMS) to encompass the organisational practices and information services through which knowledge is captured, organised, maintained, shared, and reused. Ipe's conceptual framework locates knowledge sharing in interactions between individuals, while Yeboah distinguishes its enablers, processes, and outcomes~\cite{ref095,ref093}. These perspectives explain why a searchable collection alone is insufficient: useful experience must first become accessible to others, and its reuse must serve an organisational objective. For incident response, that objective is to help people reach correct, well-supported decisions with less effort. Faster production of an answer is valuable only when its interpretation and use remain appropriate.

\section{The incident-response knowledge lifecycle}
Knowledge capture concerns both formal guidance and experience gained during investigations. NIST recommends documented procedures, including playbooks that organise actionable tasks for particular situations, and the incorporation of lessons learned into improved practice~\cite{nist2025incident}. A useful incident lesson records more than the final action: it explains the conditions under which the action was chosen and what evidence justified it. Otherwise, later responders may copy an outcome without understanding its constraints. The distinction between recording an incident and producing reusable guidance follows from the knowledge-sharing emphasis on creating, sharing, and applying knowledge~\cite{ref095}.

Organisation makes those conditions discoverable. Osman et al. compare knowledge-management tools through dimensions including knowledge sources, extraction, input, retrieval, and sharing, placing search within a broader process~\cite{ref094}. Applied to incident response, this perspective suggests describing knowledge by its purpose and context: affected technology, incident type, intended audience, author or responsible team, and date of applicability. A vendor advisory and an internal recovery procedure may concern the same vulnerability while answering different questions. Relevance therefore includes what a document enables the responder to do, not only whether it discusses the same topic.

Shared vocabularies can reduce ambiguity when knowledge crosses organisational boundaries. Takahashi and Kadobayashi's reference ontology coordinates cybersecurity operational information and existing information schemas~\cite{ref074}. An ontology defines concepts and their relationships; a knowledge graph represents particular entities and relations using such a vocabulary. These approaches can connect an attack technique to affected assets or mitigations, but structured representation does not establish that a relationship is correct. Li et al. identify noise and conflict during graph construction as sources of misleading inference~\cite{ref076}. For knowledge reuse, consistent terminology and reliable source support matter before more elaborate reasoning becomes useful.

Maintenance addresses changes in both the environment and the knowledge itself. Dhingra et al. demonstrate the difficulty of representing time-sensitive facts in language models trained on data snapshots~\cite{ref060}. The corresponding KMS problem is broader than model freshness: a formerly valid procedure can become unsuitable after a platform change, and an incident hypothesis can be corrected after further investigation. Current operational advice and a historical explanation may therefore require different versions of a source. Recording validity periods, review responsibility, and superseded guidance allows those uses to remain distinguishable rather than treating the newest text as appropriate for every question.

Sharing also depends on trust and permitted disclosure. Abraham et al.'s interview-informed perspective on cyber threat intelligence (CTI) stresses human interpretation, incompatible platforms, and different meanings of trust~\cite{ref098}. A relevant internal lesson may contain information that cannot be shared with another team or organisation. Knowledge quality, topical relevance, and permission to use the material are separate considerations. The lifecycle closes when responders can report an incorrect or incomplete recommendation and responsible contributors can revise the underlying knowledge. This feedback connects retrieval to organisational learning instead of allowing errors to recur in later answers.

\section{Machine learning for retrieval and generation}
\subsection{Learned representations and model adaptation}
Machine learning can support knowledge access without requiring a new incident classifier. A classifier assigns an input to a defined category, such as malicious or benign, whereas a representation model produces features used to compare inputs. A large language model (LLM) used for generation produces text conditioned on its input. These roles can coexist, and the task differs from the method used to train the model. Sentence-BERT, for example, learns representations through fine-tuning but subsequently uses the resulting embeddings for similarity comparisons~\cite{reimers2019sentencebert}.

An embedding is a numerical representation of text. Sentence-BERT adapts a pretrained language model so that sentences can be encoded separately and compared using measures such as cosine similarity. This permits document representations to be computed before a query arrives, avoiding a separate joint model evaluation for every possible sentence pair~\cite{reimers2019sentencebert}. In an incident-response setting, semantic retrieval may connect a description of suspicious account activity with guidance phrased in different terminology. Similarity remains a retrieval signal, however: it does not establish that the guidance is current, authorised, or correct for the affected system.

Pretraining learns parameters from earlier data; inference uses those parameters to process new inputs. Using a pretrained embedding model is therefore an application of machine learning even when no local training occurs. Fine-tuning updates model parameters for a task or domain. Retrieval instead supplies selected external material at inference time, and few-shot prompting supplies examples in the input without itself changing the parameters. These mechanisms are compatible. Lewis et al.'s original retrieval-augmented generation approach combines pretrained retrieval and generation with task-specific fine-tuning; retrieval should not be defined as the absence of training~\cite{lewis2020rag}.

\subsection{Retrieval and source-grounded generation}
Lexical and semantic retrieval offer different routes to relevant knowledge. Exact terms matter for identifiers, product names, and error messages, while learned representations can recover material with related meaning. Hybrid retrieval combines these signals. A reranker then reorders an initial candidate set using a further relevance assessment. RAGIntel illustrates this sequence through keyword and dense retrieval followed by reranking and context reduction~\cite{ref033}. The attraction for incident response is practical: precise technical references and loosely described symptoms can occur in the same query. Whether combining retrieval methods improves a particular task remains an empirical question.

Retrieval-augmented generation (RAG) conditions a language model on external information selected for the current request. Lewis et al. distinguish knowledge encoded in model parameters from knowledge held in a retrievable collection. Their model combines a neural retriever with a sequence-to-sequence generator and evaluates the combination on knowledge-intensive language tasks~\cite{lewis2020rag}. This separation provides a route for introducing revised knowledge through the collection rather than relying entirely on model retraining. For a KMS, the intended benefit is an explanation grounded in applicable sources, while the original documents remain available for inspection.

The unit of retrieval affects what evidence reaches the generator. Splitting a document into passages, often called \emph{chunks}, permits focused retrieval but can separate an instruction from its preconditions. Retrieving an entire document preserves more context at the cost of additional irrelevant material. RAGAs treats context relevance and answer faithfulness as separate concerns because finding useful information and using it correctly are different problems~\cite{ref039}. A citation consequently needs more than a matching document title: the cited passage must support the associated statement, including any qualifications that change its operational meaning.

\section{Related cybersecurity systems}
MoRSE addresses broad cybersecurity question answering through specialised retrievers over structured and unstructured knowledge. Its evaluation covers 600 questions concerning general cybersecurity, multi-document reasoning, and vulnerabilities, with expert-validated reference answers and model-assisted assessment. The study reports advantages over its selected language-model comparators~\cite{ref001}. Its contribution is evidence that organising heterogeneous cybersecurity sources around a query can improve answers in that benchmark. The evaluation unit remains the answer, however, rather than an analyst's completed incident-response task. The result therefore supports investigation of retrieval mechanisms without establishing how much time a responder would save.

GenDFIR concentrates on incident timeline analysis. It converts event data into a structured textual knowledge base, retrieves relevant events, and uses a language model without task-specific fine-tuning to produce an interpreted timeline. Its evaluation uses synthetic cyber incidents and task-specific assessments of evidence retrieval and resulting analysis~\cite{ref011}. Unlike a reference-document assistant, it retrieves observations belonging to the investigation itself. This makes event coverage and chronological relationships central to the task. Controlled scenarios allow the expected evidence to be specified, but performance on them leaves uncertainty about incomplete, inconsistent, or unfamiliar evidence in operational investigations.

RAGIntel uses MITRE ATT\&CK knowledge for attack-investigation questions and reports performance comparable to standalone language models in its evaluation~\cite{ref033}. This is a useful qualification to MoRSE's positive comparison: a retrieval stage can provide an explicit knowledge source without necessarily improving aggregate answer scores. Differences in source coverage, question type, and generator make the studies informative comparisons of approaches, not a ranking of systems under common conditions.

ForensicLLM investigates local model adaptation for digital-forensics questions. It fine-tunes a model with reduced-precision weights using question--answer material derived from research articles and curated digital artefacts. The study compares the resulting configuration with a base model and a retrieval-augmented baseline, reporting stronger correctness and relevance alongside professional assessments of the responses~\cite{ref030}. This comparison shifts attention from the retrieval mechanism alone to the complete model configuration and its domain preparation. Its local deployment also raises a distinct consideration: the feasibility of using specialist assistance where forensic information cannot be sent to an external generation service.

\begin{table}[htbp]
\centering\small
\caption{Knowledge sources and evaluation focus in selected cybersecurity studies.}
\label{tab:literature-systems}
\begin{tabularx}{\textwidth}{@{}>{\raggedright\arraybackslash}p{0.19\textwidth}>{\raggedright\arraybackslash}X>{\raggedright\arraybackslash}X@{}}
\toprule
Study & Knowledge used & Main evaluation focus \\
\midrule
MoRSE~\cite{ref001} & Heterogeneous cybersecurity references. & Answer quality across general, multi-document, and vulnerability questions. \\
GenDFIR~\cite{ref011} & Events from the incident under investigation. & Evidence retrieval and timeline analysis in synthetic incidents. \\
RAGIntel~\cite{ref033} & Structured MITRE ATT\&CK intelligence. & Responses to attack-investigation questions compared with standalone models. \\
ForensicLLM~\cite{ref030} & Forensic research and curated artefacts used for domain adaptation. & Question answering, source attribution, and professional response assessments. \\
\bottomrule
\end{tabularx}
\end{table}

Evidence from malicious-package detection further limits general claims about RAG. Ibiyo et al. find that their tested knowledge bases do not improve classification performance, while their fine-tuned condition performs substantially better~\cite{ref023}. That task predicts whether code is malicious; it does not ask whether an explanation helps an analyst use response guidance. The different evaluation targets in Table~\ref{tab:literature-systems} also explain why baseline selection must follow the decision being supported. Model choice, adaptation, retrieval, and access to additional knowledge are separate variables whose contributions cannot be inferred from the label ``RAG'' alone.

\section{Reliability, confidentiality, and adversarial content}
Retrieval introduces evidence into generation, but it also introduces an attack surface. Greshake et al. demonstrate indirect prompt injection, in which attacker-controlled content encountered by an LLM-integrated application contains instructions intended to redirect its behaviour. Their experiments include attempts to manipulate outputs and disclose information~\cite{ref045}. A cybersecurity report can therefore be useful evidence while also carrying hostile instructions. Source content needs to be treated as material to assess, not as authority to redefine the task or grant access. Separating instructions from evidence is a necessary defensive objective, but formatting or prompt wording alone does not establish resistance to these attacks.

Ordinary knowledge errors create a related reliability problem. An answer can reproduce an outdated source faithfully, omit a procedural exception, or express an uncertain incident hypothesis as fact. The Retrieval-Augmented Generation Benchmark (RGB) tests noise robustness, rejection when evidence is insufficient, integration of information, and robustness to false information; the evaluated models show particular difficulties with the latter three abilities~\cite{ref038}. For incident guidance, uncertainty and conflicting evidence therefore deserve explicit treatment rather than being hidden behind a single fluent recommendation.

Confidentiality concerns the complete information flow. Restricting document access is ineffective if restricted content can still reach an unauthorised answer, and a local generator does not by itself determine where embedding, retrieval, or evaluation services process data. Combining the trust concerns identified in CTI sharing with the disclosure risks demonstrated by prompt injection suggests treating authorisation as a constraint on eligible evidence, not a preference inferred by a language model~\cite{ref098,ref045}. Source traceability then operates within that constraint: an explanation should expose the evidence a responder is permitted to inspect without disclosing protected material.

\section{Evaluating knowledge access and response decisions}
A KMS evaluation needs to distinguish finding knowledge, explaining it, and using it in a decision. At the retrieval level, precision measures the proportion of returned material judged relevant, while recall measures how much of the annotated relevant material was recovered. Rank matters because evidence placed far down a result list may be costly to locate. Passage-level assessment also differs from document-level relevance: a document can concern the right incident type while the selected passage omits the required procedure. For access-controlled collections, the relevant set must be defined for the requesting identity and the temporal scope of the question.

RAGAs separates context relevance, answer relevance, and faithfulness. Its faithfulness method decomposes a response into statements and uses a language model to assess whether the context supports them~\cite{ref039}. ARES trains lightweight judges for related dimensions and combines their predictions with human annotations through prediction-powered inference to estimate system performance with confidence intervals~\cite{ref040}. These methods support scalable comparison, but their agreement with expert judgements needs checking in the target domain. Faithfulness to an incorrect source is not factual correctness, and factual correctness alone does not establish that a proposed action is suitable for a particular environment. Citation support and operational applicability consequently require separate examination.

Negative cases help reveal failures obscured by answerable questions. Following the distinctions examined in RGB, an assessment can include irrelevant passages, contradictory evidence, and requests for which the accessible collection contains no answer~\cite{ref038}. Appropriate abstention means recognising insufficient evidence; an unnecessary refusal when relevant evidence is available instead indicates a failure of retrieval or generation. Temporal cases add another distinction: a historical account may correctly cite superseded guidance, whereas a request for current response advice requires currently applicable knowledge. These cases test the knowledge lifecycle as well as the model's language behaviour.

For human outcomes, response quality and task efficiency need joint measurement. A search-only condition can establish how analysts perform when they inspect sources directly. Generation without retrieval tests the contribution of model knowledge, while generation with retrieval tests the additional value of contextual evidence. Comparisons are interpretable when the cases, available knowledge, and model conditions are controlled. Elapsed time should cover the response task rather than only model latency. Independently assessed decision correctness, severity of errors, and source consultation help distinguish useful acceleration from faster acceptance of unsupported advice. Usability and confidence capture the responder's experience but do not replace those outcome measures.

The organisational perspective extends evaluation beyond an individual answer. Yeboah's separation of knowledge-sharing processes and outcomes suggests examining whether new lessons are contributed and subsequently reused~\cite{ref093}. Relevant KMS measures include the time required to make approved knowledge available, whether corrections displace obsolete advice, and whether contributors can identify material needing review. A system that answers a fixed benchmark well may still lose value as its collection changes. Evaluation therefore benefits from both controlled retrieval tasks and observation of how knowledge remains usable through contribution, revision, and reuse.

\section{Synthesis and research gap}
The reviewed studies connect incident-response knowledge to several complementary research problems. Knowledge-sharing research explains why capture and maintenance depend on people and organisational purpose. Learned representations support access across differences in wording, while RAG offers a way to condition explanations on retrievable sources. Cybersecurity evaluations establish task-specific benefits and limitations, with no uniform advantage for retrieval across the compared settings~\cite{ref095,ref001,ref023,ref033}.

The remaining question is how these mechanisms support the complete path from maintained knowledge to a defensible incident-response decision. We therefore frame RQ1 around the knowledge and lifecycle needed for reuse, RQ2 around relevant retrieval and traceable guidance, and RQ3 around retrieval quality, response quality, task efficiency, and usability under confidentiality constraints. This framing treats optimisation as reduced effort accompanied by sounder decisions, not as output speed alone. Chapter~\ref{ch:system} develops the design from these requirements.
